
\subsection{Compensación del factor de potencia}

A continuación, se analizó la configuración correspondiente al esquemático visto en la Figura~\ref{fig:tp3_rlc_compensado}. Para ello, se utilizó la bobina $L_1$ y se mantuvieron cerradas las llaves $S_R$, $S_L$ y $S_C$, de modo que la carga conectada estuvo formada por una rama resistiva, una rama inductiva y un banco de capacitores conectado en paralelo.


\begin{figure}[H]
    \centering
    \begin{circuitikz}
        \draw
        (0,0) to[sV,l=$V_S$] (0,4)
              -- (7,4);

        \draw
        (0,0) -- (7,0);

        % Rama resistiva
        \draw
        (2,4) to[closing switch,l=$S_R$] (2,2.8)
              to[american resistor,l=$R$] (2,0);

        % Rama inductiva
        \draw
        (4,4) to[closing switch,l=$S_L$] (4,3)
              to[american resistor,l=$R_L$] (4,1.8)
              to[L,l=$L$] (4,0);

        % Rama capacitiva
        \draw
        (6,4) to[closing switch,l=$S_C$] (6,2.8)
              to[C,l=$C$] (6,0);
    \end{circuitikz}
    \caption{Configuración correspondiente al circuito RL compensado con un capacitor en paralelo.}
    \label{fig:tp3_rlc_compensado}
\end{figure}

\subsubsection{Corroboración experimental de los resultados}

Inicialmente, se calculó el valor teórico de capacitancia que debía colocarse en paralelo para llevar el factor de potencia lo más cerca posible de la unidad.Para esto se utilizo la siguiente ecuacion
\begin{equation}
C_{teo} = \frac{Q}{2\pi f|\underline{V}|^2 }
\label{eq:potencia_activa}
\end{equation}
Luego de utilizar los valores en la Tabla~\ref{tab:tp3_paso2} se llego al resultado $C_{\text{teo}}=\SI{7,05}{\micro\farad}$.

Luego, como el banco de capacitores solo permitía ciertos valores discretos de capacidad, se buscó experimentalmente la combinación que produjera la menor corriente total. Esto se debe a que, al cambiar solamente la capacitancia, una reducción de la corriente implica una disminución de \( |Q| \). La combinación elegida fue $C_{\text{real}}=\SI{8,0}{\micro\farad}$.


A partir de las mediciones realizadas se obtuvieron los valores presentados en las Tablas~\ref{tab:tp3_compensacion_real} y~\ref{tab:tp3_compensacion_teorica}. En esta configuración se midieron la tensión eficaz aplicada sobre la carga, la corriente eficaz total consumida por el circuito y la potencia activa entregada a la carga. A partir de estas mediciones, se calcularon la potencia aparente, la potencia reactiva, el factor de potencia y el módulo de la impedancia equivalente.

\begin{table}[H]
    \centering
    \caption{Mediciones y magnitudes calculadas para la compensación con $C_{\text{real}}$.}
    \label{tab:tp3_compensacion_real}
    \begin{tabular}{lc}
        \toprule
        Variable & Valor \\
        \midrule
        $R$ [$\Omega$] & 166,7 \\
        $R_L$ [$\Omega$] & 65,4 \\
        $C_{\text{real}}$ [$\mu$F] & 8,0 \\
        $|\underline{V}|$ [V] & 100 \\
        $|\underline{I}|$ [A] & 0,657 \\
        $P$ [W] & 68 \\
        $S$ [VA] & 68 \\
        $Q$ [VAr] & $\approx 0$ \\
        $fp$ & $\approx 1,00$ \\
        $|Z_{\text{eq}}|$ [$\Omega$] & 152,2 \\
        \bottomrule
    \end{tabular}
\end{table}

\begin{table}[H]
    \centering
    \caption{Magnitudes calculadas para la referencia con $C_{\text{teo}}$.}
    \label{tab:tp3_compensacion_teorica}
    \begin{tabular}{lc}
        \toprule
        Variable & Valor \\
        \midrule
        $R$ [$\Omega$] & 166,7 \\
        $R_L$ [$\Omega$] & 65,4 \\
        $C_{\text{teo}}$ [$\mu$F] & 7,05 \\
        $|\underline{V}|$ [V] & 100 \\
        $|\underline{I}|$ [A] & 0,66 \\
        $P$ [W] & 68 \\
        $S$ [VA] & 68 \\
        $Q$ [VAr] & $\approx 0$ \\
        $fp$ & $\approx 1,00$ \\
        $|Z_{\text{eq}}|$ [$\Omega$] & 151,5 \\
        \bottomrule
    \end{tabular}
\end{table}
\subsubsection{Triángulo de potencia}

A partir de los valores de la Tabla~\ref{tab:tp3_compensacion_real}, se construyó el triángulo de potencia correspondiente a la configuración compensada. Como se observa en la Figura~\ref{fig:triangulo_compensacion}, la potencia reactiva resultó aproximadamente nula, por lo que la potencia aparente coincidió con la potencia activa. Esto indicó que la compensación capacitiva permitió llevar el factor de potencia a un valor aproximadamente unitario.

\begin{figure}[H]
    \centering
    \triangulopot{68}{0}{68}{$R\parallel L_1\parallel C$}
    \caption{Triángulo de potencia correspondiente al circuito compensado con $C_{\text{real}}$.}
    \label{fig:triangulo_compensacion}
\end{figure}

\subsubsection{Análisis de los resultados}

Al conectar el capacitor, la corriente disminuyó de \(0{,}7\,A\) a \(0{,}657\,A\), lo cual indica que la compensación redujo la componente reactiva inductiva del circuito. Sin embargo, al calcular el factor de potencia con las mediciones experimentales se obtuvo un valor de \(1{,}035\), que no es físicamente posible. Esto se debe a que la potencia activa medida por el vatímetro fue ligeramente mayor que la potencia aparente.

La diferencia entre ambos valores fue de aproximadamente \(14{,}3\%\). Si bien el valor teórico calculado fue cercano a \(7{,}05\,\mu F\), experimentalmente se observó que la menor corriente total se obtenía con una capacitancia de \(8\,\mu F\), aun cuando el banco de capacitores permitía seleccionar un valor cercano a \(7\,\mu F\). Las razones de este error pueden ser varias. En primer lugar, el valor teórico de capacitancia fue calculado a partir de mediciones previas de tensión, corriente y potencia, por lo que cualquier error en la lectura del voltímetro, amperímetro o vatímetro se propaga al cálculo de la potencia reactiva y, en consecuencia, al valor de \(C\) necesario para compensarla. Además, la búsqueda experimental de la corriente mínima dependió de la lectura visual del amperímetro, por lo que pequeñas diferencias entre configuraciones cercanas pudieron quedar afectadas por la resolución del instrumento. Por otro lado, el modelo teórico asumía componentes ideales. Sin embargo, tanto la bobina como el banco de capacitores presentan pérdidas y comportamientos no ideales. Por estos motivos, el valor de capacitancia que minimizó la corriente experimentalmente no coincidió exactamente con el calculado de forma teórica.

Vale la pena analizar los valores medidos de vatímetro. Para la configuración compensada se obtuvo una potencia activa de \(68\,W\), y para la configuración sin capacitor se obtuvo \(66{,}4\,W\). Sin embargo, teóricamente, al cambiar únicamente la \(C\) del capacitor, no debería modificarse la potencia activa. Este error del \(2{,}35\%\) se puede atribuir a varias razones. Primero, a que el capacitor, en la realidad, se comporta como si tuviese una pequeña resistencia en paralelo, por lo que no es puramente capacitivo. Esta resistencia asociada al capacitor genera un aumento en la potencia activa. Segundo, a errores de lectura, precisión instrumental, variaciones en la tensión de alimentación y al comportamiento no ideal de los componentes.
